
\subsection{构造与收敛}
\label{subsec:6-3-3}

现在设 $(\widehat C,\delta)$ 为 d.g.c.f.z.，
且 $\widehat C^k$ 是 $C^k$ 的完备化。
我们假设 $C^k$ 在
$\Lambda_{0,\mathrm{nov}}^{(0)}$ 上自由。
令
\[
\overline C
=C/\Lambda_{0,\mathrm{nov}}^{+,(0)}C
\cong
\widehat C/
\Lambda_{0,\mathrm{nov}}^{+,(0)}\widehat C,
\]
并以 $\overline\delta$ 表示
$\overline C$ 上的诱导微分。\translatornote{原书实体页 380
把 $\overline\delta$ 称为 derivation，
又在下句使用 differential graded subalgebra；
本小节只给出了分次模与微分，
并无乘法结构。
这里按型别写为“诱导微分”与“子复形”，
见疑误 A-0367。}
我们仍将 $\overline C\subset C\subset\widehat C$
嵌入为能量零部分。
一般而言，$\overline C$
并不是 $\widehat C$ 的微分分次子复形。
设 $\{\mathbf e_i\}$ 是 $C$
在 $\Lambda_{0,\mathrm{nov}}^{(0)}$ 上的一组基，
$\{\overline{\mathbf e}_i\}$
是它在 $\overline C$
中对应的 $R$-基。
写
\[
\overline\delta(\overline{\mathbf e}_i)
=\sum_j\delta_{0,ij}\overline{\mathbf e}_j,
\]
并据此定义 $\delta_0:C\to C$：
\[
\delta_0\mathbf e_i
=\sum_j\delta_{0,ij}\mathbf e_j.
\]
并将其连续延拓至 $\widehat C$。

\begin{statement}{条件}{6.3.16}\label{cond:6-3-16}
若存在 $\lambda''>0$，
使对任意 $\lambda$ 与任意
$\mathbf v\in F^\lambda\widehat C$ 都有
\[
\delta\mathbf v-\delta_0\mathbf v
\in F^{\lambda+\lambda''}\widehat C,
\]
则称 $(\widehat C,\delta)$
满足\emph{有隙条件}。\liftIndex{I224}
\end{statement}

\begin{remark}{6.3.17}\label{rem:6-3-17}
若滤过 $A_\infty$ 代数 $(C,\mathfrak m)$
满足 $\mathfrak m_0=0$，
且依定义~\ref{def:3-2-26}
在 $G$ 上有隙，
则 $(C,\mathfrak m_1)$
满足上述有隙条件。
\end{remark}

以下假设条件~\ref{cond:6-3-16} 成立。
取常数 $\lambda_0$ 满足
$0<\lambda_0<\lambda''$，
并定义滤过
\begin{equation}
\mathfrak F^n\widehat C
=F^{n\lambda_0}\widehat C.
\tag{6.3.18}\label{eq:6-3-18}
\end{equation}
我们用这一滤过定义谱序列。
从滤过环上的滤过模构造谱序列
并非十分标准，
故为读者方便，
在此给出具体构造。
令
\begin{align}
Z_r^{p,q}(\widehat C)
&=\bigl\{x\in\mathfrak F^q\widehat C^p
\mid\delta x\in
\mathfrak F^{q+r-1}\widehat C^{p+1}\bigr\}
+\mathfrak F^{q+1}\widehat C^p,
\tag{6.3.19.1}\label{eq:6-3-19-1}\\
B_r^{p,q}(\widehat C)
&=\bigl(
\delta(\mathfrak F^{q-r+2}\widehat C^{p-1})
\cap\mathfrak F^q\widehat C^p
\bigr)+\mathfrak F^{q+1}\widehat C^p,
\tag{6.3.19.2}\label{eq:6-3-19-2}\\
E_r^{p,q}(\widehat C)
&=Z_r^{p,q}(\widehat C)
/B_r^{p,q}(\widehat C).
\tag{6.3.19.3}\label{eq:6-3-19-3}
\end{align}

记
\[
\Lambda^{(0)}(\lambda)
=\Lambda_{0,\mathrm{nov}}^{(0)}
/F^\lambda\Lambda_{0,\mathrm{nov}}^{(0)}.
\]
在 $\Lambda_{0,\mathrm{nov}}^{(0)}$
上定义滤过
$\mathfrak F^n\Lambda_{0,\mathrm{nov}}^{(0)}
=F^{n\lambda_0}\Lambda_{0,\mathrm{nov}}^{(0)}$。
其伴随分次模为
\[
\operatorname{gr}_*
(\mathfrak F\Lambda_{0,\mathrm{nov}}^{(0)})
=\bigoplus_{n\in\mathbb Z_{\ge0}}
\operatorname{gr}_n
(\mathfrak F\Lambda_{0,\mathrm{nov}}^{(0)}),
\]
其中每个
$\operatorname{gr}_n
(\mathfrak F\Lambda_{0,\mathrm{nov}}^{(0)})$
都自然同构于 $\Lambda^{(0)}(\lambda_0)$。
又有
\[
\operatorname{gr}_*
(\mathfrak F\Lambda_{0,\mathrm{nov}})
=\operatorname{gr}_*
(\mathfrak F\Lambda_{0,\mathrm{nov}}^{(0)})
[e,e^{-1}]
=\bigoplus_{n\in\mathbb Z_{\ge0}}
\operatorname{gr}_n
(\mathfrak F\Lambda_{0,\mathrm{nov}}^{(0)})
[e,e^{-1}].
\]

回忆 $E_r^{p,q}$ 自然是
$\Lambda^{(0)}(\lambda_0)$-模。
乘以 $e^{\pm1}\in\Lambda_{0,\mathrm{nov}}$
给出映射
\[
e^{\pm1}:E_r^{p,q}(\widehat C)
\longrightarrow E_r^{p\pm2,q}(\widehat C),
\]
从而使
$E_r(\widehat C)
:=\bigoplus_{p,q}E_r^{p,q}(\widehat C)$
成为
$\operatorname{gr}_*(\mathfrak F\Lambda_{0,\mathrm{nov}})$-模，
而
\[
\bigoplus_{p,q}\delta_r^{p,q}:
\bigoplus_{p,q}E_r^{p,q}(\widehat C)
\longrightarrow
\bigoplus_{p,q}E_r^{p+1,q+r-1}(\widehat C)
\]
是 $\operatorname{gr}_*
(\mathfrak F\Lambda_{0,\mathrm{nov}})$-模同态。

\begin{lemma}{6.3.20}\label{lem:6-3-20}
存在 $\Lambda^{(0)}(\lambda_0)$-模同态
\[
\delta_r^{p,q}:E_r^{p,q}(\widehat C)
\longrightarrow
E_r^{p+1,q+r-1}(\widehat C),
\]
使
\begin{align}
\delta_r^{p+1,q+r-1}\circ\delta_r^{p,q}
&=0,
\tag{6.3.21.1}\label{eq:6-3-21-1}\\
\ker\delta_r^{p,q}
/\operatorname{Im}\delta_r^{p-1,q-r+1}
&\cong E_{r+1}^{p,q}(\widehat C),
\tag{6.3.21.2}\label{eq:6-3-21-2}\\
e^{\pm1}\circ\delta_r^{p,q}
&=\delta_r^{p\pm2,q}\circ e^{\pm1}.
\tag{6.3.21.3}\label{eq:6-3-21-3}
\end{align}
\end{lemma}

\begin{proof}
定义
$\delta_r^{p,q}[x]=[\delta x]
\in E_r^{p+1,q+r-1}(\widehat C)$。
(6.3.21) 的证明是标准的。
\end{proof}

$E_r^{p,q}(\widehat C)$ 的构造当然相当标准。
与通常情形的一个区别是，
我们的滤过不是有界的：
对大的 $n$，未必有
$\mathfrak F^n\widehat C=0$。
因此，谱序列的收敛性质一般而言
远非显然。
不过它在下方稳定，
即 $\mathfrak F^0\widehat C=\widehat C$。
由此得到：

\begin{lemma}{6.3.22}\label{lem:6-3-22}
若 $q-r+2\le0$，
则存在单射
\[
E_{r+1}^{p,q}(\widehat C)
\longrightarrow E_r^{p,q}(\widehat C).
\]
\end{lemma}

\begin{proof}
当 $q-r+2\le0$ 时，
定义 $E_r^{p,q}(\widehat C)$
的分母
\[
\delta(\mathfrak F^{q-r+2}
\widehat C^{p-1})
\cap\mathfrak F^q\widehat C^p
+\mathfrak F^{q+1}\widehat C^p
\]
与 $r$ 无关。
\end{proof}

\begin{definition}{6.3.23}\label{def:6-3-23}
定义
\[
E_\infty^{p,q}(\widehat C)
=\varprojlim_r E_r^{p,q}(\widehat C)
\]
为引理~\ref{lem:6-3-22}
保证存在的\emph{逆极限}。
\end{definition}

\pagebreak
\begin{lemma}{6.3.24}\label{lem:6-3-24}
作为
$\operatorname{gr}_*(\mathfrak F\Lambda_{0,\mathrm{nov}})$-模，
有同构
\[
E_2^{*,*}(\widehat C)
\cong H(\overline C;\overline\delta)
\otimes_R
\operatorname{gr}_*(\mathfrak F\Lambda_{0,\mathrm{nov}}).
\]
\end{lemma}

\begin{proof}
依定义有
\[
E_1^{*,*}(\widehat C)
\cong\overline C\otimes_R
\operatorname{gr}_*(\mathfrak F\Lambda_{0,\mathrm{nov}}).
\]
由条件~\ref{cond:6-3-16}，
$\delta_1=\overline\delta$。
故引理得证。
\end{proof}

\begin{definition}{6.3.25}\label{def:6-3-25}
定义 $\mathfrak F^qH(\widehat C,\delta)$
为 $H(\mathfrak F^q\widehat C,\delta)$
在 $H(\widehat C,\delta)$ 中的像。
\end{definition}

为了把谱序列的极限
$E_\infty^{p,q}$
与 $\mathfrak F^qH(\widehat C,\delta)$ 联系起来，
需要下面所述的某种有限性假设。
设 $(C,\delta)$ 与 $(C',\delta')$
都是满足有隙条件的 d.g.c.f.z.。
取映射 $\varphi:C\to C'$，
满足 $\varphi\delta=\delta'\varphi$，
并设
$\overline\varphi:\overline C\to\overline C'$
为其在商模
\[
\overline C
=C/\Lambda_{0,\mathrm{nov}}^{+,(0)}C,
\qquad
\overline C'
=C'/\Lambda_{0,\mathrm{nov}}^{+,(0)}C'
\]
上的诱导映射。
利用 $C$、$C'$ 的基，
可将 $\overline\varphi$
提升为 $\varphi_0:C\to C'$。

\begin{definition}{6.3.26}\label{def:6-3-26}
在上述情形下，
若存在 $\lambda''>0$，使
\[
(\varphi-\varphi_0)
(F^\lambda\widehat C)
\subseteq F^{\lambda+\lambda''}\widehat C',
\]
则称 $\varphi:C\to C'$
满足\emph{有隙条件}，
或称其为\emph{有隙上链映射}。\liftIndex{I225}
\translatornote{原书实体页 383
在包含式右端漏写靶模的撇号，
又未明言 $\lambda''>0$。
后者与条件 6.3.16 的正能量隙同义；
此处依映射型别与该条件补明，
见疑误 A-0368--A-0369。}
\end{definition}

如注~\ref{rem:6-3-17}，
若一个滤过 $A_\infty$ 同态
依定义~\ref{def:3-2-29}
在 $G$ 上有隙，
则其诱导映射满足上述有隙条件。

\begin{definition}{6.3.27}\label{def:6-3-27}
若存在有限 d.g.c.f.z.
$(C',\delta')$，
以及有隙滤过上链映射
$\varphi:(\widehat C,\delta)
\to(C',\delta')$，
使其在上同调上诱导同构
\[
\overline\varphi_*:
H(\overline C,\overline\delta)
\cong H(\overline C',\overline\delta'),
\]
则称 $(\widehat C,\delta)$
具有\emph{弱有限性}\liftIndex{I380}。
\end{definition}

利用这些定义，我们证明：

\begin{theorem}{6.3.28}\label{thm:6-3-28}
\liftIndex{I369}
若 $C$ 弱有限，
则存在 $r_0$，使下列同构成立。\translatornote{原书实体页 383
把 $r_0$ 写成对全体 $p,q$ 统一的页数；
下文证明中的常数 $c$ 却可随次数 $p$ 变化。
按定义 6.3.8，“有限”不要求次数有界；
现有假设只支持逐固定次数的稳定界，
见疑误 A-0370。}
\[
E_{r_0}^{p,q}(\widehat C)
\cong E_{r_0+1}^{p,q}(\widehat C)
\cong\cdots\cong
E_\infty^{p,q}(\widehat C)
\cong
\mathfrak F^qH^p(\widehat C,\delta)
/\mathfrak F^{q+1}H^p(\widehat C,\delta)
\]
它们都是
$\Lambda^{(0)}(\lambda_0)
=\Lambda_{0,\mathrm{nov}}^{(0)}
/F^{\lambda_0}\Lambda_{0,\mathrm{nov}}^{(0)}$-模。
\end{theorem}

\begin{proof}
先考虑 $C^k$ 为有限 c.f.z. 的情形。
此时 $C^k=\widehat C^k$。
由命题~\ref{prop:6-3-9}，
存在正数 $c$，使对任意 $\lambda$ 都有
\[
\delta(C^k)\cap F^\lambda C^{k+1}
\subseteq\delta(F^{\lambda-c}C^k).
\]
取 $r_0$ 满足
$(r_0-1)\lambda_0-c>\lambda_0$。
现令 $r\ge r_0$，
$\mathbf x\in Z_r^{p,q}(C)$。
可写 $\mathbf x=\mathbf x_0+\mathbf x_1$，
其中
\[
\delta\mathbf x_0
\in\mathfrak F^{r+q-1}C
=F^{(r+q-1)\lambda_0}C,
\qquad
\mathbf x_1\in\mathfrak F^{q+1}C.
\]
因此由命题~\ref{prop:6-3-9}，
存在 $\mathbf y$，使
\[
\delta\mathbf x_0=\delta\mathbf y,
\qquad
\mathbf y\in
F^{(r+q-1)\lambda_0-c}C
\subset F^{(q+1)\lambda_0}C
=\mathfrak F^{q+1}C.
\]
于是
\[
[\mathbf x]
=[\mathbf x-\mathbf y]
=[\mathbf x_0-\mathbf y]
\in E_r^{p,q}(C),
\]
并且
\[
\delta_r[\mathbf x]
=\delta_r[\mathbf x_0-\mathbf y]=0,
\qquad
E_{r_0}^{p,q}\cong E_{r_0+1}^{p,q}
\cong\cdots\cong E_\infty^{p,q}.
\]

为证
$E_\infty^{p,q}
\cong\mathfrak F^qH^p(C,\delta)
/\mathfrak F^{q+1}H^p(C,\delta)$，
先构造映射
\[
\pi_{p,q}:
\mathfrak F^qH^p(C,\delta)
\longrightarrow E_\infty^{p,q}.
\]
设 $\mathbf x\in\mathfrak F^q C^p$
且 $\delta\mathbf x=0$。
则对任意 $r\ge r_0$，
$[\mathbf x]\in E_r^{p,q}(C)$；
这些类给出 $E_\infty^{p,q}(C)$
中的元素。
若 $\mathbf x-\mathbf x'=\delta\boldsymbol\eta$，
取 $r>q+2$，则
$\boldsymbol\eta\in
\mathfrak F^{q-r+2}C^{p-1}=C^{p-1}$，
故 $[\mathbf x]=[\mathbf x']$
于 $E_r^{p,q}(C)$。\translatornote{原书实体页 383
此处把当前复形 $C$ 误写成旧记号 $V$，
且只取 $r>q$；
但 (6.3.19.2) 的分母要求
$\boldsymbol\eta\in
\mathfrak F^{q-r+2}C^{p-1}$。
取 $r>q+2$ 即足够，
见疑误 A-0371--A-0372。}
因此此赋值给出良定义的 $\pi_{p,q}$。

下面证明 $\pi_{p,q}$ 满射。
取
$[\mathbf x]\in E_{r_0}^{p,q}(C)
\cong E_\infty^{p,q}(C)$。
由前面的论证，
存在 $\mathbf y\in\mathfrak F^{q+1}C$，使
\[
[\mathbf x-\mathbf y]
=[\mathbf x]\in E_r^{p,q}(C),
\qquad
\delta(\mathbf x-\mathbf y)=0.
\]
于是 $\mathbf x-\mathbf y$
给出 $\mathfrak F^qC^p$
中的上闭链，
而其上同调类在 $\pi_{p,q}$ 下
映到 $[\mathbf x]$。
故 $\pi_{p,q}$ 满射。

再证
$\ker\pi_{p,q}
=\mathfrak F^{q+1}H^p(C,\delta)$。
取 $[\mathbf x]\in
\mathfrak F^qH^p(C,\delta)$，
且 $\pi_{p,q}([\mathbf x])=0$。
当 $r>q+2$ 时，有
\[
\mathbf x\in
\delta(\mathfrak F^{q-r+2}C^{p-1})
+\mathfrak F^{q+1}C^p
\subseteq
\delta(C^{p-1})+\mathfrak F^{q+1}C^p.
\]
因此
$[\mathbf x]\in\mathfrak F^{q+1}H^p(C,\delta)$，
得到
$\ker\pi_{p,q}\subseteq
\mathfrak F^{q+1}H^p(C,\delta)$。
反向包含显然成立。

最后撤去 $\widehat C$ 有限生成的条件。
依定义~\ref{def:6-3-27}，
取有限 d.g.c.f.z.
$(C',\delta')$，
以及诱导约化上同调同构的
有隙上链映射
$\varphi:\widehat C\to C'$。
于是 $\varphi$ 诱导谱序列同态。
它在 $E_2$ 项是同构，
故在 $E_\infty$ 项也是同构。
另一方面，$\overline\varphi$
依假设诱导
$H(\overline C,\overline\delta)
\to H(\overline C',\overline\delta')$
的同构。
仿照命题~\ref{prop:4-3-18} 的证明，
可证 $\varphi$ 诱导上同调同构
\[
H(\widehat C,\delta)
\longrightarrow H(C',\delta').
\]
因此，既然对 $C'$ 而言
$\pi_{p,q}$ 是同构，
对 $\widehat C$ 亦然。
定理~\ref{thm:6-3-28} 得证。
\end{proof}

对周期哈密顿系统的弗洛尔同调，
利用与命题~\ref{prop:6-3-9}
类似的断言，
可证明 \textbf{[Ono95]} 中
的修正弗洛尔同调
与通常的弗洛尔同调同构。
